\section{Classical baseline: thresholding and morphology}
\label{sec:classical}

Before any learned method, we establish what classical image processing achieves on the same protocol: a deep model is only interesting if it beats interpretable, millisecond-scale processing. The \emph{crater detector} exploits the fact that craters appear as locally dark depressions under the fixed WAC illumination: each tile is normalised, adaptive (local-mean) thresholding isolates dark regions, and morphological opening/closing plus a minimum-area filter clean the result into a binary crater mask. The pipeline has five interpretable parameters (threshold offset, window size, two kernel radii, minimum area), tuned on the calibration subset of the spatial validation split only. Wrinkle ridges are curvilinear, so the \emph{ridge detector} instead uses the Sato vesselness filter (Hessian eigenvalues at multiple scales), which responds to elongated bright structures, followed by the same morphological cleaning and threshold calibration.

The baseline runs at ${\sim}1$\,ms per tile on a single CPU core --- three orders of magnitude cheaper than Mask R-CNN --- and every decision it makes can be inspected. Its scores in Sect.~\ref{sec:comparison} define the floor that justifies (or fails to justify) the learned methods class by class.
